% ==============================================================================
% פרק 6: היתכנות טכנולוגית, מפת דרכים ליישום ומטריצת סיכונים
% ==============================================================================
\chapter{היתכנות טכנולוגית, מפת דרכים ליישום ומטריצת סיכונים}
\label{chap:implementation_roadmap}

\section{מחסנית טכנולוגית מוצעת ובחירת סביבות עבודה}
\label{sec:impl_tech_stack_selection}
היתכנות המימוש של פלטפורמת \projectname{} נשענת על ההתקדמות המהירה של השנים האחרונות בשלושה תחומים מרכזיים: מזעור מודלי שפה לעיבוד מקומי בקצה (\textenglish{Small Language Models - SLMs}), מסגרות תזמור מרובות סוכנים מבוססות גרפים, ומסדי נתונים וקטוריים קלי-משקל.

מחסנית הטכנולוגיה המוצעת לפיתוח המערכת כוללת ארבע שכבות ליבה:
\begin{enumerate}
    \item \textbf{מודלי שפה מקומיים לעיבוד בקצה (\textenglish{Edge AI Inference})} --
    כדי לעמוד בדרישת הברזל של שמירה מוחלטת על פרטיות המשתמש, המערכת תתבסס על מודלים בעלי ארכיטקטורה קומפקטית ומכומתת (\textenglish{4-bit Quantized GGUF}), כגון \textenglish{Llama-3-8B-Instruct}~\cite{touvron2023llama} או \textenglish{Phi-3-Mini} (3.8 מיליארד פרמטרים). מודלים אלו מסוגלים לרוץ מקומית על גבי מעבד גרפי של מחשב נייד או שבב \textenglish{NPU} מודרני בסמארטפונים בקצב של 25--40 טוקנים לשנייה, תוך צריכת זיכרון (\textenglish{VRAM}) הנמוכה מ־6 ג'יגה-בייט.
    \item \textbf{מסגרת תזמור מרובת סוכנים (\textenglish{Multi-Agent Orchestration Framework})} --
    אימוץ סביבת \textenglish{LangGraph} או \textenglish{AutoGen}~\cite{wu2023autogen}. סביבות אלו מאפשרות תיאום א-סינכרוני מבוסס גרף מצבים מחזורי (\textenglish{Cyclic State Graph}), ניהול זיכרון משותף בין סוכנים, טיפול בשגיאות תקשורת, ועקיבות מלאה של נתיבי ההחלטה באמצעות רישום אירועים (\textenglish{Event Sourcing}).
    \item \textbf{מסד נתונים וקטורי מקומי (\textenglish{Local Vector Database})} --
    שימוש בספריית \textenglish{ChromaDB} או \textenglish{SQLite-vec} לשמירת היסטוריית החשיפות ורפלקציות העבר של סוכן הרפלקציה. המערכת תבצע חיפושי דמיון קוסינוס מקומיים ללא כל פנייה לשרתי ענן, תוך שימוש במודל שיכון קומפקטי מקומי (כגון \textenglish{all-MiniLM-L6-v2}).
    \item \textbf{שכבת ממשק משתמש ואיסוף נתונים (\textenglish{HCI Frontend \& Data Ingestion})} --
    תוסף דפדפן חוצה-פלטפורמות (\textenglish{Chrome/Firefox Manifest V3 Extension}) או אפליקציית דסקטופ מבוססת \textenglish{Electron/React}, המאפשרת ממשק חזותי עשיר, הצגת עקומות המצב הנפשי בזמן אמת, והענקת פיגומים מטא-קוגניטיביים עדינים.
\end{enumerate}

\section{מבנה פירוק עבודה (WBS), אבני דרך ותרשים גאנט}
\label{sec:impl_milestones_gantt}
תהליך הפיתוח המתוכנן של פרויקט \projectname{} מאורגן בחמש חבילות עבודה מובנות (\textenglish{Work Packages - WPs}), הנפרסות על פני מחזור פיתוח שלם בן 14 שבועות סמסטריאליים (או לחלופין 12 חודשי מחקר מורחבים):

\begin{enumerate}
    \item \textbf{\textenglish{WP1}: מפרט מושגי, תיקוף מתודולוגי וארכיטקטורה}
    \begin{itemize}
        \item \textbf{משימות}: \textenglish{T1.1} אפיון שאלות המחקר ומדדי היעד הכמותיים; \textenglish{T1.2} אימות ספרות אמפירית וכיול משוואות וקטור המצב $\wellbeingVector(t)$; \textenglish{T1.3} הגדרת סכמות \textenglish{JSON-RPC 2.0} ופרוטוקול התקשורת.
        \item \textbf{תוצר מסירה (\textenglish{Deliverable D1})}: דוח מכין אקדמי שלם (מטלה 2) ומפרט דרישות הנדסי מתוקף.
    \end{itemize}

    \item \textbf{\textenglish{WP2}: תשתית קלט מקומית וסיווג סמנטי מרובה-סוכנים}
    \begin{itemize}
        \item \textbf{משימות}: \textenglish{T2.1} פיתוח תוסף דפדפן ומנגנון אנונימיזציה מקומי (\agentObserver{}); \textenglish{T2.2} כתיבת פרומפטים מערכתיים וכיול מדדי $\comparisonScore$ ו־$\exposureIndex$ (\agentAnalyst{}); \textenglish{T2.3} תיקוף דיוק סיווג ($F_1 \ge 0.85$) מול קורפוס מוזהב של 200 פוסטים מתויגים.
        \item \textbf{תוצר מסירה (\textenglish{Deliverable D2})}: אב-טיפוס עובד של מודול קליטה וסיווג מקומי עם דוח ביצועים.
    \end{itemize}

    \item \textbf{\textenglish{WP3}: מנוע רפלקציה, זיכרון אפיזודי ומיקרו-התערבויות}
    \begin{itemize}
        \item \textbf{משימות}: \textenglish{T3.1} הטמעת מסד נתונים וקטורי מקומי (\textenglish{ChromaDB}) ופונקציית שליפה תלת-ממדית; \textenglish{T3.2} מימוש משוואות דעיכה ועדכון וקטור המצב $\wellbeingVector(t)$ (\agentReflector{}); \textenglish{T3.3} כיול מנגנון הניסוח הסוקרטי ומניעת ריאקטנס (\agentInterventionist{}).
        \item \textbf{תוצר מסירה (\textenglish{Deliverable D3})}: ליבת מערכת \textenglish{MAS} מתואמת הפועלת מקומית בקצה.
    \end{itemize}

    \item \textbf{\textenglish{WP4}: מערך סימולציה מבוקר עם 20 פרסונות}
    \begin{itemize}
        \item \textbf{משימות}: \textenglish{T4.1} יצירת 20 פרסונות גנרטיביות על פני 4 ארכיטיפים פסיכולוגיים; \textenglish{T4.2} הרצת 1,800 ריצות סימולציה מבוקרות (30 חזרות $\times$ 3 תנאים); \textenglish{T4.3} חישוב שיעור בלימת שחיקה ($\text{DMR} \ge 40\%$) וניתוח מתאם מול שאלוני \textenglish{RSES} ו־\textenglish{PANAS}.
        \item \textbf{תוצר מסירה (\textenglish{Deliverable D4})}: מסד נתוני סימולציה מלא ודוח אימות מתודולוגי מקיף.
    \end{itemize}

    \item \textbf{\textenglish{WP5}: ניסוי פיילוט משתמשים, לוח מחוונים ודוח מסכם}
    \begin{itemize}
        \item \textbf{משימות}: \textenglish{T5.1} הגשת בקשה וקבלת אישור ועדת אתיקה (\textenglish{IRB}) באוניברסיטת חיפה; \textenglish{T5.2} ניסוי פיילוט בקרב 50 סטודנטים ומדידת שביעות רצון ואי-חודרנות; \textenglish{T5.3} ניתוח סטטיסטי מסכם וכתיבת דוח גמר אקדמי.
        \item \textbf{תוצר מסירה (\textenglish{Deliverable D5})}: פלטפורמת "מראה קוגניטיבית" פעילה ודוח מחקר מסכם.
    \end{itemize}
\end{enumerate}

\subsection{מטריצת אחריות תפעולית עתידית (RACI)}
\label{subsec:impl_raci_matrix}

\begin{agentbox}{הערה מתודולוגית מנחה אודות מטריצת ה־RACI}
חלוקת תפקידי ה־\textenglish{RACI} שלהלן מוצעת כמודל תפעולי ליישום עתידי בלבד, והיא כפופה לאישור שתי המגישות. הקרדיט על כתיבת הדוח המכין הנוכחי, המחקר האקדמי והגדרת הפרויקט הינו שוויוני ומשותף לחלוטין בין שמס סבאג ומלאק פרוג'ה.
\end{agentbox}

טבלה~\ref{tab:raci_matrix} מפרטת את הקצאת התפקידים המוצעת לחבילות העבודה: אחראי ביצוע (\textenglish{Responsible - R}), בעל אחריות כוללת (\textenglish{Accountable - A}), גורם נועץ (\textenglish{Consulted - C}), וגורם מעודכן (\textenglish{Informed - I}):

\begin{table}[H]
\centering
\caption{מטריצת אחריות תפעולית מוצעת ליישום עתידי (RACI Matrix)}
\label{tab:raci_matrix}
\small
\begin{tabularx}{\textwidth}{>{\bfseries\color{DarkPurple}\raggedleft\arraybackslash}p{3.4cm} >{\centering\arraybackslash}p{2.6cm} >{\centering\arraybackslash}p{2.6cm} >{\centering\arraybackslash}p{2.6cm} >{\centering\arraybackslash}p{2.2cm}}
\toprule
\rowcolor{LavenderCard}
\textbf{\color{DarkPurple}חבילת עבודה / משימה} & \textbf{\color{DarkPurple}שמס סבאג \newline {\tiny(פסיכולוגיה ומתודולוגיה)}} & \textbf{\color{DarkPurple}מלאק פרוג'ה \newline {\tiny(תקשורת, משפיענים וממשק)}} & \textbf{\color{DarkPurple}אחריות כוללת (A)} & \textbf{\color{DarkPurple}גורם נועץ (C) / מעודכן (I)} \\
\midrule
\textenglish{WP1}: מפרט ודוח מכין & \textenglish{R} & \textenglish{R} & משותף (\textenglish{A}) & מרצה הקורס (\textenglish{C}) \\
\midrule
\textenglish{WP2}: קליטה וסיווג משפיענים & \textenglish{C} & \textenglish{R} & מלאק (\textenglish{A}) & מרצה הקורס (\textenglish{I}) \\
\midrule
\textenglish{WP3}: מנוע רפלקציה ופסיכומטריה & \textenglish{R} & \textenglish{C} & שמס (\textenglish{A}) & מרצה הקורס (\textenglish{I}) \\
\midrule
\textenglish{WP4}: מערך סימולציה ופרסונות & \textenglish{R} & \textenglish{R} & משותף (\textenglish{A}) & מרצה הקורס (\textenglish{C}) \\
\midrule
\textenglish{WP5}: פיילוט משתמשים ודוח סופי & \textenglish{R} & \textenglish{R} & משותף (\textenglish{A}) & ועדת אתיקה / מרצה (\textenglish{C/A}) \\
\bottomrule
\end{tabularx}
\end{table}

\subsection{תרשים גאנט, אבני דרך והנתיב הקריטי}
\label{subsec:impl_gantt_critical_path}
איור~\ref{fig:project_gantt_chart} מציג את תרשים הגאנט הכרונולוגי של חמש חבילות העבודה, אבני הדרך ותוצרי המסירה (\textenglish{D1--D5}), תוך הדגשת \textbf{הנתיב הקריטי} (\textenglish{Critical Path}) המסומן בקו כפול אדום.

\begin{figure}[H]
\centering
\begin{LTR}
\begin{tikzpicture}[
    scale=0.82, every node/.style={transform shape},
    bar/.style={rectangle, rounded corners=1.5mm, line width=0.8pt, inner sep=4pt, font=\scriptsize\sffamily},
    timeline/.style={draw=LavenderBorder, line width=0.8pt},
    milestone/.style={diamond, draw=DarkPurple, fill=PurpleAccent, inner sep=2.5pt}
]

% Time Axis Grid (Weeks 1 to 14)
\foreach \w in {1,...,14} {
    \draw[timeline] (\w*0.95, 0) -- (\w*0.95, 5.8);
    \node[font=\tiny\bfseries\color{SlateMuted}, below=0.1cm of {(\w*0.95, 0)}] {W\w};
}
\draw[line width=1.1pt, color=DarkPurple] (0.95, 0) -- (14.25, 0);

% Phase 1 / WP1 (W1-W3)
\filldraw[fill=DarkPurple, draw=DarkPurple] (0.95, 4.8) rectangle (2.85, 5.3);
\node[anchor=west, font=\scriptsize\bfseries\color{DarkPurple}] at (3.0, 5.05) {WP1: Specification \& Modeling};
\node[milestone] at (2.85, 5.05) {};
\node[above=2pt, font=\tiny\bfseries\color{DarkPurple}] at (2.85, 5.3) {D1};

% Phase 2 / WP2 (W3-W6)
\filldraw[fill=RoyalPurple, draw=RoyalPurple] (2.85, 3.8) rectangle (5.70, 4.3);
\node[anchor=west, font=\scriptsize\bfseries\color{RoyalPurple}] at (5.85, 4.05) {WP2: Ingestion \& Classifier Prototyping};
\node[milestone] at (5.70, 4.05) {};
\node[above=2pt, font=\tiny\bfseries\color{DarkPurple}] at (5.70, 4.3) {D2};

% Phase 3 / WP3 (W6-W9)
\filldraw[fill=PurpleAccent, draw=DarkPurple] (5.70, 2.8) rectangle (8.55, 3.3);
\node[anchor=west, font=\scriptsize\bfseries\color{DarkPurple}] at (8.70, 3.05) {WP3: Reflector \& Interventionist Synthesis};
\node[milestone] at (8.55, 3.05) {};
\node[above=2pt, font=\tiny\bfseries\color{DarkPurple}] at (8.55, 3.3) {D3};

% Phase 4 / WP4 (W9-W11)
\filldraw[fill=AlertRose, draw=AlertRose] (8.55, 1.8) rectangle (10.45, 2.3);
\node[anchor=west, font=\scriptsize\bfseries\color{AlertRose}] at (10.60, 2.05) {WP4: 20-Persona Simulation Experiments};
\node[milestone] at (10.45, 2.05) {};
\node[above=2pt, font=\tiny\bfseries\color{DarkPurple}] at (10.45, 2.3) {D4};

% Phase 5 / WP5 (W11-W14)
\filldraw[fill=SuccessEmerald, draw=SuccessEmerald] (10.45, 0.8) rectangle (13.30, 1.3);
\node[anchor=west, font=\scriptsize\bfseries\color{SuccessEmerald}] at (13.45, 1.05) {WP5: Pilot Trials \& Final Report};
\node[milestone] at (13.30, 1.05) {};
\node[above=2pt, font=\tiny\bfseries\color{DarkPurple}] at (13.30, 1.3) {D5};

% Critical Path Dependencies (Red Highlighted Line)
\draw[->, line width=1.5pt, color=BrickRed] (2.85, 4.8) -- (2.85, 4.3);
\draw[->, line width=1.5pt, color=BrickRed] (5.70, 3.8) -- (5.70, 3.3);
\draw[->, line width=1.5pt, color=BrickRed] (8.55, 2.8) -- (8.55, 2.3);
\draw[->, line width=1.5pt, color=BrickRed] (10.45, 1.8) -- (10.45, 1.3);

% Legend
\node[anchor=north west, font=\tiny\sffamily, align=left, fill=white, draw=LavenderBorder, rounded corners=1.5mm, inner sep=4pt] at (0.95, 2.5) {
\textbf{Legend:}\\
{\color{BrickRed}\textbf{=== Critical Path (WP1 $\to$ WP5)}}\\
{\color{DarkPurple}$\blacklozenge$ Deliverable Milestone (D1--D5)}\\
\textit{* 14-week schedule subject to Haifa lecturer confirmation}
};

\end{tikzpicture}
\end{LTR}
\caption{תרשים גאנט לפרויקט \projectname{}: חבילות עבודה (WP1--WP5), תוצרי מסירה (D1--D5) והנתיב הקריטי (הערה: לוח הזמנים בן 14 שבועות סמסטריאליים כפוף לאישור המרצה באוניברסיטת חיפה)}
\label{fig:project_gantt_chart}
\end{figure}

\section{מטריצת סיכונים כמותית ($5 \times 5$) ואסטרטגיות הפחתה}
\label{sec:impl_risk_matrix}
ניהול סיכונים בפרויקט \projectname{} מושתת על מתודולוגיית הערכה כמותית במטריצת $5 \times 5$. ציון הסיכון הכולל ($R$) מחושב כמכפלת סבירות ההתרחשות ($P \in [1, 5]$) בחומרת ההשפעה ($I \in [1, 5]$):
\begin{equation}
    R = P \times I \quad (1 \le R \le 25)
\end{equation}

הסולם הכמותי מוגדר כדלקמן:
\begin{itemize}
    \item \textbf{סבירות ($P$)}: 1 = נדירה מאוד ($<5\%$), 2 = נמוכה ($5\text{--}20\%$), 3 = בינונית ($20\text{--}50\%$), 4 = גבוהה ($50\text{--}80\%$), 5 = כמעט ודאית ($>80\%$).
    \item \textbf{השפעה ($I$)}: 1 = זניחה, 2 = מתונה, 3 = משמעותית, 4 = חמורה, 5 = קריטית להמשך הפרויקט.
    \item \textbf{סף סיכון קריטי}: כל סיכון שציונו $R \ge 12$ מוגדר כ\textbf{סיכון קריטי} ומסומן באדום, ומחייב תוכנית מניעה מוקדמת ותוכנית מגירה אופרטיבית מיידית.
    \item \textbf{ציון שיורי ($R_{\text{res}}$)}: ציון הסיכון המחושב לאחר יישום פעולות המניעה והבקרה ($R_{\text{res}} = P_{\text{res}} \times I_{\text{res}}$).
\end{itemize}

טבלה~\ref{tab:risk_matrix} מפרטת את חמשת סיכוני הליבה של הפרויקט, את ציוני הסיכון הגולמיים והשיוריים, ואת תוכניות המניעה והמגירה:

\begin{table}[H]
\centering
\caption{מטריצת סיכונים כמותית ($5 \times 5$), פעולות מניעה, תוכניות מגירה וציונים שיוריים}
\label{tab:risk_matrix}
\scriptsize
\begin{tabularx}{\textwidth}{>{\bfseries\color{DarkPurple}}p{1.1cm} >{\raggedleft\arraybackslash}p{2.3cm} >{\centering\arraybackslash}p{0.7cm} >{\centering\arraybackslash}p{0.7cm} >{\centering\arraybackslash}p{0.7cm} >{\centering\arraybackslash}p{1.1cm} X >{\raggedleft\arraybackslash}p{2.3cm} >{\centering\arraybackslash}p{0.8cm}}
\toprule
\rowcolor{LavenderCard}
\textbf{\color{DarkPurple}מזהה} & \textbf{\color{DarkPurple}תיאור הסיכון} & \textbf{\color{DarkPurple}$P$} & \textbf{\color{DarkPurple}$I$} & \textbf{\color{DarkPurple}$R$} & \textbf{\color{DarkPurple}רמת סיכון} & \textbf{\color{DarkPurple}פעולת מניעה והפחתה} & \textbf{\color{DarkPurple}תוכנית מגירה} & \textbf{\color{DarkPurple}$R_{\text{res}}$} \\
\midrule
\textenglish{RSK-01} & הזיות מודל שפה בסיווג פסיכולוגי & 3 & 4 & \textbf{\color{BrickRed}12} & \textbf{\color{BrickRed}קריטי} & כפיית פלט \textenglish{JSON Schema}, רובריקות סיווג חדות וטמפרטורה $\tau = 0.2$. & כיול מול מאגר 200 פוסטים ושימוש במודל שיפוט (\textenglish{Judge LLM}). & 3 \newline {\tiny($1 \times 3$)} \\
\midrule
\textenglish{RSK-02} & שיהוי ועומס חישובי בקצה & 4 & 3 & \textbf{\color{BrickRed}12} & \textbf{\color{BrickRed}קריטי} & מודלי \textenglish{4-bit GGUF}, עיבוד באצוות (\textenglish{Batching}) וסינון יוריסטי מוקדם. & העברת עיבוד כבד למודל שפה מבוזר מקומי ברשת הביתית. & 4 \newline {\tiny($2 \times 2$)} \\
\midrule
\textenglish{RSK-03} & התנגדות פסיכולוגית (\textenglish{Reactance}) & 3 & 4 & \textbf{\color{BrickRed}12} & \textbf{\color{BrickRed}קריטי} & שפה סוקרטית עדינה, מגבלת 2 התראות/שעה ושליטה מלאה של המשתמש. & מעבר לתצוגה שקטה בלוח המחוונים ללא התראות קופצות. & 3 \newline {\tiny($1 \times 3$)} \\
\midrule
\textenglish{RSK-04} & שבירת קליטת \textenglish{DOM} ברשתות & 4 & 3 & \textbf{\color{BrickRed}12} & \textbf{\color{BrickRed}קריטי} & שימוש בסלקטורים מבוססי מטא-נתונים כלליים ובידוד שכבת הקלט לתוסף. & עדכון דינמי של חוקי החילוץ באמצעות קובץ קונפיגורציה מקומי. & 4 \newline {\tiny($2 \times 2$)} \\
\midrule
\textenglish{RSK-05} & עיכוב באישור ועדת אתיקה (\textenglish{IRB}) & 2 & 4 & 8 & בינוני & הגשת בקשה מוקדמת מבוססת פרוטוקול נתונים מקומי ללא איסוף מידע מזהה. & הרחבת מערך הסימולציה עם 40 פרסונות כחלופה עד לקבלת האישור. & 3 \newline {\tiny($1 \times 3$)} \\
\bottomrule
\end{tabularx}
\end{table}

\section{הקצאת משאבים והיתכנות חישובית באוניברסיטת חיפה}
\label{sec:impl_resource_allocation}
ההיתכנות החישובית של פרויקט \projectname{} נשענת על ניצול יעיל של משאבי החומרה הקיימים במעבדות המחקר של אוניברסיטת חיפה.

עבור הרצת מערך הסימולציה עם 20 הפרסונות הגנרטיביות, נדרש שרת מחשוב בעל מעבד גרפי מדגם \textenglish{NVIDIA RTX 4090} (24 ג'יגה-בייט \textenglish{VRAM}) או שרת ענן אקדמי ייעודי. משאב זה מאפשר הרצה מקבילית של 4 מופעי סוכנים מקומיים בו-זמנית, תוך השלמת סשן סימולציה מלא של 100 פוסטים בתוך פחות מ־45 דקות.

עלויות הפיתוח הכוללות נשמרות ברמה מינימלית בזכות השימוש במודלי קוד פתוח (\textenglish{Open-Weight Models}) ובכלי תוכנה חופשיים, עובדה המבטיחה את שחזוריות המחקר והתאמתו למסגרות אקדמיות עתידיות.
